\textbf{Phần 1. Trắc nghiệm lựa chọn}
\setcounterpageref{ex}{0}

\begin{ex} %[3D4N23]
Kết quả của phép nhân $12 \times 3$ là
\choice
{$32$}
{$38$}
{\True $36$}
{$40$}
\loigiai{
$12 \times 3 = 36$.\\
\textbf{Đáp án của câu hỏi là C}
}
\end{ex}

\begin{ex} %[3D4N23]
Kết quả của phép nhân $20 \times 4$ là
\choice
{$60$}
{$100$}
{$70$}
{\True $80$}
\loigiai{
$20 \times 4 = 80$.\\
\textbf{Đáp án của câu hỏi là D}
}
\end{ex}

\begin{ex} %[3D4H23]
Kết quả của phép nhân có nhớ $14 \times 5$ là
\choice
{$60$}
{\True $70$}
{$80$}
{$65$}
\loigiai{
$14 \times 5 = 70$.\\
\textbf{Đáp án của câu hỏi là B}
}
\end{ex}

\begin{ex} %[3D4H23]
Kết quả của phép nhân $26 \times 3$ là
\choice
{\True $78$}
{$68$}
{$88$}
{$72$}
\loigiai{
$26 \times 3 = 78$.\\
\textbf{Đáp án của câu hỏi là A}
}
\end{ex}

\begin{ex} %[3D4H23]
Tính giá trị biểu thức: $15 \times 4 + 20$.
\choice
{$70$}
{\True $80$}
{$90$}
{$85$}
\loigiai{
$15 \times 4 + 20 = 60 + 20 = 80$.\\
\textbf{Đáp án của câu hỏi là B}
}
\end{ex}

\begin{ex} %[3D4H23]
Tìm $x$ biết: $x : 4 = 18$.
\choice
{$68$}
{$76$}
{\True $72$}
{$80$}
\loigiai{
$x = 18 \times 4 = 72$.\\
\textbf{Đáp án của câu hỏi là C}
}
\end{ex}

\begin{ex} %[3D4V23]
Mỗi gói kẹo có 24 cái kẹo. Hỏi 3 gói kẹo như thế có tất cả bao nhiêu cái kẹo?
\choice
{$68$ cái}
{$76$ cái}
{$64$ cái}
{\True $72$ cái}
\loigiai{
$24 \times 3 = 72$ cái kẹo.\\
\textbf{Đáp án của câu hỏi là D}
}
\end{ex}

\begin{ex} %[3D4V23]
Một cuộn vải may được 16 bộ quần áo. Hỏi 5 cuộn vải như thế may được bao nhiêu bộ quần áo?
\choice
{$75$ bộ}
{\True $80$ bộ}
{$85$ bộ}
{$90$ bộ}
\loigiai{
$16 \times 5 = 80$ bộ quần áo.\\
\textbf{Đáp án của câu hỏi là B}
}
\end{ex}

\begin{ex} %[3D4V23]
Một trang sách có 32 dòng chữ, mỗi dòng có khoảng 3 từ. Hỏi trang sách đó có khoảng bao nhiêu từ?
\choice
{\True $96$ từ}
{$86$ từ}
{$92$ từ}
{$106$ từ}
\loigiai{
$32 \times 3 = 96$ từ.\\
\textbf{Đáp án của câu hỏi là A}
}
\end{ex}

\begin{ex} %[3D4V23]
Tính giá trị biểu thức: $13 \times 6 - 28$.
\choice
{$48$}
{$52$}
{$54$}
{\True $50$}
\loigiai{
$13 \times 6 - 28 = 78 - 28 = 50$.\\
\textbf{Đáp án của câu hỏi là D}
}
\end{ex}

\begin{ex} %[3D4V23]
Lớp 3A có 3 tổ, mỗi tổ có 12 học sinh. Tổng số học sinh của lớp 3A là
\choice
{$32$ học sinh}
{$38$ học sinh}
{\True $36$ học sinh}
{$40$ học sinh}
\loigiai{
$12 \times 3 = 36$ học sinh.\\
\textbf{Đáp án của câu hỏi là C}
}
\end{ex}

\begin{ex} %[3D4C23]
Một hộp có 15 viên bi đỏ và 18 viên bi xanh. Hỏi 3 hộp như thế có tất cả bao nhiêu viên bi?
\choice
{\True $99$ viên}
{$96$ viên}
{$90$ viên}
{$102$ viên}
\loigiai{
Mỗi hộp có: $15 + 18 = 33$ viên bi.\\
3 hộp có: $33 \times 3 = 99$ viên bi.\\
\textbf{Đáp án của câu hỏi là A}
}
\end{ex}

\textbf{Phần 2. Trắc nghiệm đúng sai}
\setcounterpageref{ex}{0}

\begin{ex} %[3D4H23]
Xét tính đúng sai của các phép nhân:
\choiceTF
{\True $13 \times 3 = 39$.}
{\True $16 \times 4 = 64$.}
{$25 \times 3 = 65$.}
{\True $25 \times 3 = 75$.}
\loigiai{
\begin{itemize}
    \item $13 \times 3 = 39$: Đúng.
    \item $16 \times 4 = 64$: Đúng.
    \item $25 \times 3 = 75 \neq 65$. Ý 3 Sai, Ý 4 Đúng.
\end{itemize}
\begin{itemize}
    \item a) \textbf{Đúng}.
    \item b) \textbf{Đúng}.
    \item c) \textbf{Sai}.
    \item d) \textbf{Đúng}.
\end{itemize}
}
\end{ex}

\begin{ex} %[3D4V23]
Một người nuôi 4 lồng chim, mỗi lồng có 14 con chim. Xét tính đúng sai:
\choiceTF
{\True Cả 4 lồng có tất cả 56 con chim.}
{\True Nếu bán đi 20 con chim thì còn lại 36 con.}
{Nếu mua thêm 15 con chim thì tổng cộng có 70 con.}
{\True Nếu mua thêm 15 con chim thì tổng cộng có 71 con.}
\loigiai{
\begin{itemize}
    \item Tổng: $14 \times 4 = 56$ con. Đúng.
    \item Còn lại: $56 - 20 = 36$ con. Đúng.
    \item Thêm 15 con: $56 + 15 = 71$ con $\neq 70$. Ý 3 Sai, Ý 4 Đúng.
\end{itemize}
\begin{itemize}
    \item a) \textbf{Đúng}.
    \item b) \textbf{Đúng}.
    \item c) \textbf{Sai}.
    \item d) \textbf{Đúng}.
\end{itemize}
}
\end{ex}

\textbf{Phần 3. Trả lời ngắn}
\setcounterpageref{ex}{0}

\begin{ex} %[3D4H23]
Tính tích: $23 \times 3$.
\shortans{69}
\loigiai{
\textbf{Đáp số:} $69$\\
$23 \times 3 = 69$.
}
\end{ex}

\begin{ex} %[3D4H23]
Tính tích: $18 \times 4$.
\shortans{72}
\loigiai{
\textbf{Đáp số:} $72$\\
$18 \times 4 = 72$.
}
\end{ex}

\begin{ex} %[3D4V23]
Tính giá trị biểu thức: $16 \times 5 + 15$.
\shortans{95}
\loigiai{
\textbf{Đáp số:} $95$\\
$80 + 15 = 95$.
}
\end{ex}

\begin{ex} %[3D4V23]
Một giá sách có 4 ngăn, mỗi ngăn có 22 cuốn sách. Tính tổng số cuốn sách có trên giá.
\shortans{88}
\loigiai{
\textbf{Đáp số:} $88$\\
$22 \times 4 = 88$ cuốn sách.
}
\end{ex}

\begin{ex} %[3D4V23]
Tìm $x$ biết: $x : 6 = 14$.
\shortans{84}
\loigiai{
\textbf{Đáp số:} $84$\\
$x = 14 \times 6 = 84$.
}
\end{ex}

\begin{ex} %[3D4C23]
Một hình vuông có cạnh dài $18\text{ cm}$. Tính chu vi hình vuông đó theo xăng-ti-mét.
\shortans{72}
\loigiai{
\textbf{Đáp số:} $72$\\
$18 \times 4 = 72\text{ cm}$.
}
\end{ex}
